% Part: applied-modal-logics
% Chapter: epistemic-logic
% Section: properties-accessibility

\documentclass[../../../include/open-logic-section]{subfiles}

\begin{document}

\olfileid{aml}{el}{acc}

\olsection{प्राप्यता संबंध आणि ज्ञानविषयक तत्त्वे}

सामान्य मोडल तर्कशास्त्रांतील चौकट-अनुरूपतेबद्दल आपल्याला
आधीच जे माहीत आहे, त्या आधारावर ज्ञानविषयक अर्थलक्षण
दिल्यास वैशिष्ट्यदर्शक !!{formula}s कशी दिसतात हे
पाहता येईल. ज्ञानविषयक तर्कशास्त्रांचा अर्थ सहसा S5
मध्ये लावतात, हे आपण सांगितले आहे. आता विविध
ज्ञानविषयक तत्त्वे कशी व्यक्त होतात आणि त्यांचा
विविध चौकट-अटींशी कसा अनुरूप संबंध आहे ते पाहू.

सामान्य मोडल तर्कशास्त्रात विविध मोडल !!{formula}s
प्राप्यता संबंधांचे निरनिराळे गुणधर्म वैशिष्ट्यित
करतात, हे आठवा. विशिष्ट ज्ञानविषयक तत्त्वांशी
संबंधित अशी काही सूत्रे या तक्त्यात दिली आहेत.

\begin{table}[t]
    \begin{tabular}{| p{.48\textwidth} || p{.48\textwidth} |}
      \hline
      {\emph{$R$ असा असेल \dots}} & {\emph{तर~$\mModel{M}$ मध्ये \dots सत्य असते:}} \\
      \hline \hline

      & $\Knows (p \lif q) \lif (\Knows p \lif \Knows q)$
      \hfill \newline{(संवरण)} \\
      \hline
      \emph{परावर्ती}: $\forall w Rww$
      & $\Knows p \lif p$ \hfill \newline{(ज्ञानाची सत्यता)} \\
      \hline
      \emph{संक्रमक}: \newline
      $\forall u \forall v \forall w ((Ruv \land Rvw) \lif Ruw)$ &
      $\Knows p \lif \Knows \Knows p$ \hfill
           \newline{(सकारात्मक आत्मनिरीक्षण)} \\
      \hline
      \emph{युक्लिडी}: \newline
      $\forall w \forall u \forall v ((Rwu \land Rwv) \lif Ruv)$ &
      $\lnot \Knows p \lif \Knows \neg \Knows p$ \hfill
        \newline{(नकारात्मक आत्मनिरीक्षण)} \\
      \hline
    \end{tabular}
    \caption{चार ज्ञानविषयक तत्त्वे.}
    \ollabel{tab:four}
  \end{table}

$T$ स्वयंसिद्धाशी संबंधित ज्ञानाची सत्यता हे या
तत्त्वांपैकी सर्वात कमी विवाद्य मानतात, कारण
!!a{formula} माहीत असेल, तर ते सत्य असलेच पाहिजे,
असा तिचा आशय आहे. संवरण तसेच सकारात्मक आणि
नकारात्मक आत्मनिरीक्षण ही तत्त्वे त्यापेक्षा बरीच
अधिक विवाद्य आहेत.

$K$ स्वयंसिद्धाशी संबंधित संवरण म्हणजे कर्त्याचे
ज्ञान निहितार्थाखाली बंद असते, ही कल्पना. काही
प्रसंगी ती रास्त वाटू शकते. उदाहरणार्थ, मी
व्हिक्टोरियात असेन, तर व्हॅन्कुव्हर बेटावर आहे,
हे मला माहीत असू शकते. काही विलक्षण संशयवादी
परिस्थिती बाजूला ठेवल्या, तर मी व्हिक्टोरियात
आहे हे मला माहीत आहे; म्हणून मी व्हॅन्कुव्हर
बेटावर आहे हेही मला माहीत असावे. या बाबतीत
माझ्या ज्ञानाचे तार्किक संवरण बऱ्यापैकी सहज
पटण्यासारखे वाटेल. दुसरीकडे, आपल्या ज्ञानाचे
परिणाम आपण नेहमी विचारात घेत नाही; म्हणून इतर
प्रसंगी हा निष्कर्ष तितका सहज पटणार नाही.

सकारात्मक आत्मनिरीक्षणाला कधीकधी KK-तत्त्व म्हणतात:
मला एखादी गोष्ट माहीत असेल, तर ती मला माहीत आहे
हेही मला माहीत असते, असे ते सांगते. ते 4 स्वयंसिद्धाचे
ज्ञानविषयक समरूप आहे. त्याचप्रमाणे नकारात्मक
आत्मनिरीक्षण असे सांगते की एखादी गोष्ट मला
\emph{माहीत नसेल}, तर ती मला माहीत नाही हे मला
माहीत असते; ते 5 स्वयंसिद्धाचे समरूप आहे. या
दोन्हींसाठी बऱ्यापैकी नेहमीची प्रतिउदाहरणे देता
येतात: एखादी गोष्ट मला प्रत्यक्षात माहीत असतानाही,
ती मला माहीत आहे की नाही याबद्दल मी साशंक असू
शकतो.


\end{document}
